\documentclass[11pt]{article}

\usepackage[margin=1in]{geometry}
\usepackage{amsmath,amssymb,amsthm}
\usepackage{enumitem}
\usepackage{xcolor}
\usepackage{hyperref}

\hypersetup{
  colorlinks=true,
  linkcolor=blue!60!black,
  urlcolor=blue!60!black
}

\newcommand{\R}{\mathbb{R}}
\newcommand{\N}{\mathbb{N}}
\newcommand{\B}{\mathcal{B}}
\newcommand{\F}{\mathcal{F}}
\newcommand{\A}{\mathcal{A}}
\newcommand{\C}{\mathcal{C}}
\newcommand{\Sd}{\mathcal{S}_d}
\newcommand{\Prob}{\mathbb{P}}
\newcommand{\Knowledge}[1]{%
  \par\smallskip\noindent
  \textbf{Knowledge used.}
  {\small\color{blue!60!black}\raggedright #1\par}
  \par\smallskip
}

\title{Chapter 1 Exercises}
\author{\textsc{Solutions to Rick Durrett's}\\
\textit{Probability: Theory and Examples}}
\date{\today}

\begin{document}
\maketitle

\noindent
These solutions use the local Chapter 1 LLM/Viki knowledge set in
\path{Probability/LLM_Viki_Dataset/Chapter1_Knowledge}.  The cited
knowledge IDs refer to entries in
\path{chapter1_knowledge_pieces.jsonl}.

\section*{Exercise 1.1.1}

\noindent\textbf{Question.}
Let $\Omega=\R$, let $\F$ be the collection of all subsets $A$ such that
$A$ or $A^c$ is countable, and define
\[
  P(A)=
  \begin{cases}
  0, & A \text{ is countable},\\
  1, & A^c \text{ is countable}.
  \end{cases}
\]
Show that $(\Omega,\F,P)$ is a probability space.

\Knowledge{
\path{durrett_ch1_probability_space};
\path{durrett_ch1_sigma_field};
\path{durrett_ch1_measure_continuity}.
}

\noindent\textbf{Solution.}
First we show that $\F$ is a $\sigma$-field.  Since $\varnothing$ is
countable, $\varnothing\in\F$ and $\Omega\in\F$.  If $A\in\F$, then
either $A$ is countable or $A^c$ is countable, so the same alternative
holds for $A^c$; hence $A^c\in\F$.

Now let $A_1,A_2,\ldots\in\F$.  If every $A_n$ is countable, then
$\bigcup_n A_n$ is countable, so $\bigcup_n A_n\in\F$.  If at least one
$A_k$ is uncountable, then $A_k^c$ is countable.  Since
\[
  \left(\bigcup_{n=1}^\infty A_n\right)^c
  = \bigcap_{n=1}^\infty A_n^c
  \subseteq A_k^c,
\]
the complement of the union is countable.  Therefore
$\bigcup_n A_n\in\F$.  Thus $\F$ is a $\sigma$-field.

The definition of $P$ is consistent: no set can be both countable and
co-countable in $\R$, because $\R$ is uncountable.  Also $P(\Omega)=1$.
It remains to prove countable additivity.  Let $A_1,A_2,\ldots$ be
pairwise disjoint members of $\F$.

If all $A_n$ are countable, then $\bigcup_n A_n$ is countable and
\[
  P\left(\bigcup_{n=1}^\infty A_n\right)=0
  =\sum_{n=1}^\infty P(A_n).
\]
If some $A_k$ is co-countable, then no other disjoint $A_n$ can be
co-countable.  Indeed, for $n\ne k$ we have
$A_n\subseteq A_k^c$, so $A_n$ is countable.  Hence
\[
  P\left(\bigcup_{n=1}^\infty A_n\right)=1
  =P(A_k)+\sum_{n\ne k}P(A_n).
\]
Therefore $P$ is countably additive and $(\Omega,\F,P)$ is a probability
space.

\section*{Exercise 1.1.2}

\noindent\textbf{Question.}
Recall the definition of $\Sd$ from Example 1.1.5:
\[
  \Sd=\{\varnothing\}\cup
  \left\{(a_1,b_1]\times\cdots\times(a_d,b_d]:
  -\infty\le a_i<b_i\le \infty\right\}.
\]
Show that $\sigma(\Sd)=\B(\R^d)$, the Borel subsets of $\R^d$.

\Knowledge{
\path{durrett_ch1_sigma_field};
\path{durrett_ch1_semialgebra_extension};
\path{durrett_ch1_product_distribution_function}.
}

\noindent\textbf{Solution.}
Every set in $\Sd$ is a product of half-open intervals, each of which is
Borel in $\R$.  Hence every element of $\Sd$ is Borel in $\R^d$, and
\[
  \sigma(\Sd)\subseteq \B(\R^d).
\]

For the reverse inclusion, it is enough to show that every open subset of
$\R^d$ lies in $\sigma(\Sd)$, because the Borel $\sigma$-field is
generated by the open sets.  First observe that if
$q_i<r_i$ are rational numbers, then the open rational rectangle
\[
  Q=(q_1,r_1)\times\cdots\times(q_d,r_d)
\]
belongs to $\sigma(\Sd)$.  Indeed,
\[
  Q
  =\bigcup_{m=1}^\infty
  (q_1,r_1-1/m]\times\cdots\times(q_d,r_d-1/m],
\]
where terms with $r_i-1/m\le q_i$ are simply omitted.  Each displayed
half-open rectangle belongs to $\Sd$, so $Q\in\sigma(\Sd)$.

Now let $G\subseteq\R^d$ be open.  Since $\R^d$ has a countable base of
open rectangles with rational endpoints, $G$ is a countable union of
such rational rectangles.  Each rational rectangle is in $\sigma(\Sd)$,
so $G\in\sigma(\Sd)$.  Thus $\B(\R^d)\subseteq\sigma(\Sd)$.
Combining both inclusions gives $\sigma(\Sd)=\B(\R^d)$.

\section*{Exercise 1.1.3}

\noindent\textbf{Question.}
A $\sigma$-field $\F$ is said to be countably generated if there is a
countable collection $\C\subseteq\F$ such that $\sigma(\C)=\F$.
Show that $\B(\R^d)$ is countably generated.

\Knowledge{
\path{durrett_ch1_sigma_field};
\path{durrett_ch1_semialgebra_extension};
\path{durrett_ch1_product_distribution_function}.
}

\noindent\textbf{Solution.}
Let $\C$ be the collection of all open balls in $\R^d$ whose centers have
rational coordinates and whose radii are positive rational numbers:
\[
  \C=\{B(q,r): q\in\mathbb{Q}^d,\ r\in\mathbb{Q},\ r>0\}.
\]
This collection is countable because $\mathbb{Q}^d\times\mathbb{Q}_+$ is
countable.

We claim that $\C$ is a base for the usual topology on $\R^d$.  Let
$G\subseteq\R^d$ be open and let $x\in G$.  Choose $\varepsilon>0$ such
that $B(x,\varepsilon)\subseteq G$.  Pick a rational point
$q\in\mathbb{Q}^d$ close enough to $x$ and a rational radius $r>0$ small
enough so that
\[
  x\in B(q,r)\subseteq B(x,\varepsilon)\subseteq G.
\]
Therefore every open set is a union of sets from $\C$.  Since $\C$ is
countable, every open set is in fact a countable union of members of
$\C$.

Thus all open sets belong to $\sigma(\C)$, so
$\B(\R^d)\subseteq\sigma(\C)$.  Conversely, every member of $\C$ is open
and hence Borel, so $\sigma(\C)\subseteq\B(\R^d)$.  Therefore
\[
  \sigma(\C)=\B(\R^d),
\]
and $\B(\R^d)$ is countably generated.

\section*{Exercise 1.1.4}

\noindent\textbf{Question.}
\begin{enumerate}[label=(\roman*)]
\item Show that if $\F_1\subseteq\F_2\subseteq\cdots$ are
$\sigma$-algebras, then $\bigcup_i\F_i$ is an algebra.
\item Give an example to show that $\bigcup_i\F_i$ need not be a
$\sigma$-algebra.
\end{enumerate}

\Knowledge{
\path{durrett_ch1_sigma_field}.
}

\noindent\textbf{Solution.}
Let
\[
  \F_\infty=\bigcup_{i=1}^\infty \F_i.
\]
First, $\Omega\in\F_1\subseteq\F_\infty$.  If $A\in\F_\infty$, then
$A\in\F_i$ for some $i$.  Since $\F_i$ is a $\sigma$-algebra,
$A^c\in\F_i\subseteq\F_\infty$.

Now let $A,B\in\F_\infty$.  Then $A\in\F_i$ and $B\in\F_j$ for some
$i,j$.  If $k=\max(i,j)$, then $A,B\in\F_k$ because the sequence is
increasing.  Hence $A\cup B\in\F_k\subseteq\F_\infty$.  Thus
$\F_\infty$ is closed under complements and finite unions, so it is an
algebra.

For a counterexample to countable closure, take $\Omega=\N$ and let
\[
  \F_n=\sigma(\{1\},\{2\},\ldots,\{n\}).
\]
Equivalently, $\F_n$ consists of all sets whose behavior on
$\{1,\ldots,n\}$ is arbitrary but whose tail $\{n+1,n+2,\ldots\}$ is
either entirely included or entirely excluded.  Then
\[
  \F_1\subseteq\F_2\subseteq\cdots
\]
are $\sigma$-algebras.  Their union contains every finite subset of
$\N$, because any finite set is contained in $\{1,\ldots,n\}$ for large
enough $n$.  Hence each singleton $\{2k\}$ belongs to
$\bigcup_n\F_n$.  But
\[
  E=\{2,4,6,\ldots\}
  =\bigcup_{k=1}^\infty \{2k\}
\]
is neither finite nor cofinite, so it does not belong to any $\F_n$.
Therefore $\bigcup_n\F_n$ is not closed under countable unions and is not
a $\sigma$-algebra.

\section*{Exercise 1.1.5}

\noindent\textbf{Question.}
A set $A\subseteq\{1,2,\ldots\}$ is said to have asymptotic density
$\theta$ if
\[
  \lim_{n\to\infty}
  \frac{|A\cap\{1,2,\ldots,n\}|}{n}
  =\theta.
\]
Let $\A$ be the collection of sets for which the asymptotic density
exists.  Is $\A$ a $\sigma$-algebra?  Is $\A$ an algebra?

\Knowledge{
\path{durrett_ch1_sigma_field};
\path{durrett_ch1_measure_continuity}.
}

\noindent\textbf{Solution.}
The collection $\A$ is not an algebra, and therefore it is not a
$\sigma$-algebra either.  We construct two sets with asymptotic density
but whose intersection does not have asymptotic density.

Let $N_0=0$ and choose a rapidly increasing sequence, for example
$N_k=2^{2^k}$.  Define blocks
\[
  I_k=\{N_{k-1}+1,\ldots,N_k\},\qquad k\ge 1.
\]
Let
\[
  E=\{2,4,6,\ldots\}
\]
be the even integers, so $E$ has asymptotic density $1/2$.

Now define a set $B$ block by block:
\[
  B\cap I_k=
  \begin{cases}
  E\cap I_k, & k \text{ odd},\\
  E^c\cap I_k, & k \text{ even}.
  \end{cases}
\]
In every block, $B$ contains exactly one parity class, so up to a bounded
rounding error it contains half of each block.  Consequently
\[
  \frac{|B\cap\{1,\ldots,n\}|}{n}\to \frac12,
\]
and hence $B\in\A$.

Consider now $E\cap B$.  On odd blocks, $E\cap B$ contains essentially
all even numbers in the block; on even blocks, it contains none of the
even numbers in the block.  Because $N_k/N_{k-1}\to\infty$, the most
recent block dominates all previous blocks.  Therefore, along the end of
odd blocks,
\[
  \frac{|(E\cap B)\cap\{1,\ldots,N_k\}|}{N_k}\to \frac12,
\]
whereas along the end of even blocks,
\[
  \frac{|(E\cap B)\cap\{1,\ldots,N_k\}|}{N_k}\to 0.
\]
The density of $E\cap B$ does not exist.

Thus $E\in\A$ and $B\in\A$, but $E\cap B\notin\A$.  Hence $\A$ is not
closed under finite intersections, so it is not an algebra.  Since every
$\sigma$-algebra is an algebra, $\A$ is not a $\sigma$-algebra.

\section*{Exercise 1.2.1}

\noindent\textbf{Question.}
Let $X$ and $Y$ be random variables on $(\Omega,\F,P)$ and let
$A\in\F$.  Define $Z=X$ on $A$ and $Z=Y$ on $A^c$.  Show that $Z$ is a
random variable.

\Knowledge{
\path{durrett_ch1_measurable_map};
\path{durrett_ch1_measurability_generator};
\path{durrett_ch1_sigma_field}.
}

\noindent\textbf{Solution.}
It is enough to check that $\{Z\le x\}\in\F$ for each $x\in\R$.  For any
real $x$,
\[
  \{Z\le x\}
  =
  \bigl(A\cap\{X\le x\}\bigr)
  \cup
  \bigl(A^c\cap\{Y\le x\}\bigr).
\]
Since $A\in\F$, $A^c\in\F$, and $X,Y$ are random variables, both
$\{X\le x\}$ and $\{Y\le x\}$ are in $\F$.  The right-hand side is formed
using finite intersections and a finite union, so it belongs to $\F$.
Therefore $Z$ is measurable and hence is a random variable.

\section*{Exercise 1.2.2}

\noindent\textbf{Question.}
Let $\chi$ be standard normal.  Use Durrett's normal-tail estimate to
find upper and lower bounds for $P(\chi\ge 4)$.

\Knowledge{
\path{durrett_ch1_normal_tail_bound};
\path{durrett_ch1_distribution_function}.
}

\noindent\textbf{Solution.}
For a standard normal random variable,
\[
  P(\chi\ge 4)
  =
  \frac{1}{\sqrt{2\pi}}\int_4^\infty e^{-y^2/2}\,dy.
\]
Durrett's Theorem 1.2.6 says that, for $x>0$,
\[
  (x^{-1}-x^{-3})e^{-x^2/2}
  \le
  \int_x^\infty e^{-y^2/2}\,dy
  \le
  x^{-1}e^{-x^2/2}.
\]
With $x=4$, this gives
\[
  \frac{1}{\sqrt{2\pi}}\left(\frac14-\frac{1}{64}\right)e^{-8}
  \le
  P(\chi\ge 4)
  \le
  \frac{1}{\sqrt{2\pi}}\cdot\frac14 e^{-8}.
\]
Equivalently,
\[
  \boxed{
  \frac{15}{64\sqrt{2\pi}}e^{-8}
  \le
  P(\chi\ge 4)
  \le
  \frac{1}{4\sqrt{2\pi}}e^{-8}}
\]
which is approximately
\[
  3.135\times 10^{-5}
  \le
  P(\chi\ge 4)
  \le
  3.346\times 10^{-5}.
\]

\section*{Exercise 1.2.3}

\noindent\textbf{Question.}
Show that a distribution function has at most countably many
discontinuities.

\Knowledge{
\path{durrett_ch1_distribution_function};
\path{durrett_ch1_atoms_jumps};
\path{durrett_ch1_measure_continuity}.
}

\noindent\textbf{Solution.}
Let $F$ be a distribution function.  Since $F$ is nondecreasing, the left
limit
\[
  F(x-)=\lim_{y\uparrow x}F(y)
\]
exists for every $x$.  A discontinuity of a distribution function is a
jump, so $x$ is a discontinuity exactly when
\[
  \Delta F(x):=F(x)-F(x-)>0.
\]

For integers $m,n\ge1$, define
\[
  D_{m,n}
  =
  \{x\in[-n,n]: \Delta F(x)>1/m\}.
\]
We claim $D_{m,n}$ is finite.  If
$x_1<\cdots<x_k$ are distinct points in $D_{m,n}$, then the jumps at
these points are disjoint contributions to the total increase of $F$ on
an interval containing $[-n,n]$.  Hence
\[
  \sum_{j=1}^k \Delta F(x_j)\le 1.
\]
But each jump is greater than $1/m$, so $k<m$.  Thus $D_{m,n}$ is finite.

Every discontinuity has a positive jump, so it lies in $D_{m,n}$ for
some $m,n$.  Therefore the set of discontinuities is contained in
\[
  \bigcup_{m=1}^\infty\bigcup_{n=1}^\infty D_{m,n},
\]
a countable union of finite sets.  Hence a distribution function has at
most countably many discontinuities.

\section*{Exercise 1.2.4}

\noindent\textbf{Question.}
Suppose $F(x)=P(X\le x)$ is continuous.  Show that $Y=F(X)$ is uniform
on $(0,1)$, i.e. $P(Y\le y)=y$ for $0\le y\le1$.

\Knowledge{
\path{durrett_ch1_probability_integral_transform};
\path{durrett_ch1_distribution_function};
\path{durrett_ch1_cdf_to_law}.
}

\noindent\textbf{Solution.}
Fix $y\in[0,1]$.  Define
\[
  x_y=\sup\{x\in\R:F(x)\le y\}.
\]
Using the continuity of $F$ and the limits $F(-\infty)=0$,
$F(\infty)=1$, we have $F(x_y)=y$, with the obvious interpretation at
$y=0$ and $y=1$.

Because $F$ is nondecreasing,
\[
  \{F(X)\le y\}=\{X\le x_y\}
\]
up to possible endpoints in flat parts of $F$; continuity ensures those
endpoints do not change the probability beyond the value $F(x_y)$.  More
formally, the set $\{x:F(x)\le y\}$ is a left ray ending at $x_y$,
possibly with a flat interval at level $y$, and its probability is
$F(x_y)=y$.  Therefore
\[
  P(Y\le y)=P(F(X)\le y)=P(X\le x_y)=F(x_y)=y.
\]
Thus $Y$ has the uniform distribution on $(0,1)$.

\section*{Exercise 1.2.5}

\noindent\textbf{Question.}
Suppose $X$ has continuous density $f$ and is supported on
$[\alpha,\beta]$.  Let $g$ be strictly increasing and differentiable on
$(\alpha,\beta)$.  Find the density of $g(X)$.

\Knowledge{
\path{durrett_ch1_change_of_variables};
\path{durrett_ch1_distribution_function};
\path{durrett_ch1_cdf_to_law}.
}

\noindent\textbf{Solution.}
Let $Y=g(X)$.  Since $g$ is strictly increasing, it has an inverse on
$(g(\alpha),g(\beta))$.  For $y\in(g(\alpha),g(\beta))$,
\[
  P(Y\le y)
  =
  P(g(X)\le y)
  =
  P(X\le g^{-1}(y))
  =
  \int_\alpha^{g^{-1}(y)} f(x)\,dx.
\]
Differentiating with respect to $y$ gives
\[
  f_Y(y)
  =
  f(g^{-1}(y))\frac{d}{dy}g^{-1}(y)
  =
  \frac{f(g^{-1}(y))}{g'(g^{-1}(y))},
  \qquad g(\alpha)<y<g(\beta),
\]
where the displayed formula is understood at points where
$g'(g^{-1}(y))>0$.  Outside the interval $(g(\alpha),g(\beta))$, the
density is $0$.

In particular, if $g(x)=ax+b$ with $a>0$, then $g^{-1}(y)=(y-b)/a$ and
\[
  f_Y(y)=\frac1a f\left(\frac{y-b}{a}\right).
\]

\section*{Exercise 1.2.6}

\noindent\textbf{Question.}
Suppose $X$ is normally distributed.  Use the previous exercise to find
the density of $\exp(X)$, the lognormal distribution.

\Knowledge{
\path{durrett_ch1_change_of_variables};
\path{durrett_ch1_common_distributions_expectations};
\path{durrett_ch1_distribution_function}.
}

\noindent\textbf{Solution.}
Let $X\sim N(\mu,\sigma^2)$ with $\sigma>0$, and set
\[
  Y=e^X.
\]
Here $g(x)=e^x$, so $g^{-1}(y)=\log y$ and
$g'(g^{-1}(y))=e^{\log y}=y$.  The density of $X$ is
\[
  f_X(x)=\frac{1}{\sigma\sqrt{2\pi}}
  \exp\left[-\frac{(x-\mu)^2}{2\sigma^2}\right].
\]
Therefore, for $y>0$,
\[
  f_Y(y)
  =
  \frac{f_X(\log y)}{y}
  =
  \frac{1}{y\sigma\sqrt{2\pi}}
  \exp\left[-\frac{(\log y-\mu)^2}{2\sigma^2}\right].
\]
For $y\le0$, $f_Y(y)=0$.  Thus
\[
  \boxed{
  f_Y(y)=
  \frac{1}{y\sigma\sqrt{2\pi}}
  \exp\left[-\frac{(\log y-\mu)^2}{2\sigma^2}\right]\mathbf{1}_{(0,\infty)}(y)}
\]
is the lognormal density.

\section*{Exercise 1.2.7}

\noindent\textbf{Question.}
\begin{enumerate}[label=(\roman*)]
\item If $X$ has density $f$, compute the distribution function and
density of $X^2$.
\item Specialize to $X$ standard normal to get the chi-square density
with one degree of freedom.
\end{enumerate}

\Knowledge{
\path{durrett_ch1_change_of_variables};
\path{durrett_ch1_distribution_function};
\path{durrett_ch1_common_distributions_expectations}.
}

\noindent\textbf{Solution.}
Let $Y=X^2$.  For $y<0$, clearly $P(Y\le y)=0$.  For $y\ge0$,
\[
  F_Y(y)
  =
  P(X^2\le y)
  =
  P(-\sqrt y\le X\le \sqrt y)
  =
  \int_{-\sqrt y}^{\sqrt y} f(x)\,dx.
\]
Differentiating for $y>0$ gives
\[
  f_Y(y)
  =
  f(\sqrt y)\frac{1}{2\sqrt y}
  +
  f(-\sqrt y)\frac{1}{2\sqrt y}
  =
  \frac{f(\sqrt y)+f(-\sqrt y)}{2\sqrt y},
  \qquad y>0.
\]
Thus
\[
  f_Y(y)=
  \frac{f(\sqrt y)+f(-\sqrt y)}{2\sqrt y}\mathbf{1}_{(0,\infty)}(y).
\]

If $X$ is standard normal, then
\[
  f(x)=\frac{1}{\sqrt{2\pi}}e^{-x^2/2},
\]
so $f(\sqrt y)=f(-\sqrt y)=(2\pi)^{-1/2}e^{-y/2}$.  Hence, for $y>0$,
\[
  f_Y(y)
  =
  \frac{2(2\pi)^{-1/2}e^{-y/2}}{2\sqrt y}
  =
  \frac{1}{\sqrt{2\pi y}}e^{-y/2}.
\]
Therefore $X^2$ has density
\[
  \boxed{
  f_Y(y)=
  \frac{1}{\sqrt{2\pi y}}e^{-y/2}\mathbf{1}_{(0,\infty)}(y)}
\]
which is the chi-square distribution with one degree of freedom.

\section*{Exercise 1.3.1}

\noindent\textbf{Question.}
Show that if $\A$ generates $\mathcal{S}$, then
\[
  X^{-1}(\A)\equiv \{\{X\in A\}:A\in\A\}
\]
generates
\[
  \sigma(X)=\{\{X\in B\}:B\in\mathcal{S}\}.
\]

\Knowledge{
\path{durrett_ch1_measurability_generator};
\path{durrett_ch1_sigma_field};
\path{durrett_ch1_measurable_map}.
}

\noindent\textbf{Solution.}
Let
\[
  \mathcal{G}=\sigma\bigl(X^{-1}(\A)\bigr).
\]
Since $\A\subseteq\mathcal{S}$, we immediately have
$X^{-1}(\A)\subseteq\sigma(X)$, so
\[
  \mathcal{G}\subseteq\sigma(X).
\]

For the reverse inclusion, define
\[
  \mathcal{D}=\{B\in\mathcal{S}:\{X\in B\}\in\mathcal{G}\}.
\]
The class $\mathcal{D}$ is a $\sigma$-field: inverse images commute with
complements and countable unions.  Also $\A\subseteq\mathcal{D}$ by the
definition of $\mathcal{G}$.  Since $\A$ generates $\mathcal{S}$, it
follows that $\mathcal{D}=\mathcal{S}$.  Therefore for every
$B\in\mathcal{S}$, $\{X\in B\}\in\mathcal{G}$, so
\[
  \sigma(X)\subseteq\mathcal{G}.
\]
Thus $\sigma(X)=\sigma(X^{-1}(\A))$.

\section*{Exercise 1.3.2}

\noindent\textbf{Question.}
Prove Theorem 1.3.6 when $n=2$ by checking directly that
$\{X_1+X_2<x\}\in\F$.

\Knowledge{
\path{durrett_ch1_measurable_arithmetic};
\path{durrett_ch1_measurability_generator};
\path{durrett_ch1_sigma_field}.
}

\noindent\textbf{Solution.}
Let $X_1$ and $X_2$ be random variables.  For fixed $x\in\R$, observe
that
\[
  \{X_1+X_2<x\}
  =
  \bigcup_{q\in\mathbb{Q}}
  \bigl(\{X_1<q\}\cap\{X_2<x-q\}\bigr).
\]
Indeed, if $X_1(\omega)+X_2(\omega)<x$, then there is a rational
$q$ with
\[
  X_1(\omega)<q<x-X_2(\omega).
\]
Conversely, if $X_1(\omega)<q$ and $X_2(\omega)<x-q$, then
$X_1(\omega)+X_2(\omega)<x$.

For each rational $q$, the events $\{X_1<q\}$ and $\{X_2<x-q\}$ are in
$\F$.  The rationals are countable, so the union above is a countable
union of measurable events.  Hence $\{X_1+X_2<x\}\in\F$ for every
$x\in\R$.  Therefore $X_1+X_2$ is a random variable.

\section*{Exercise 1.3.3}

\noindent\textbf{Question.}
Show that if $f$ is continuous and $X_n\to X$ almost surely, then
$f(X_n)\to f(X)$ almost surely.

\Knowledge{
\path{durrett_ch1_composition_measurable};
\path{durrett_ch1_limits_measurable}.
}

\noindent\textbf{Solution.}
Let
\[
  \Omega_0=\{\omega:X_n(\omega)\to X(\omega)\}.
\]
By assumption, $P(\Omega_0)=1$.  If $\omega\in\Omega_0$, then
$X_n(\omega)\to X(\omega)$ as real numbers.  Since $f$ is continuous,
\[
  f(X_n(\omega))\to f(X(\omega)).
\]
Thus convergence holds for every $\omega\in\Omega_0$, a probability-one
set.  Therefore $f(X_n)\to f(X)$ almost surely.

\section*{Exercise 1.3.4}

\noindent\textbf{Question.}
\begin{enumerate}[label=(\roman*)]
\item Show that a continuous function from $\R^d$ to $\R$ is measurable
from $(\R^d,\B(\R^d))$ to $(\R,\B(\R))$.
\item Show that $\B(\R^d)$ is the smallest $\sigma$-field that makes all
continuous functions measurable.
\end{enumerate}

\Knowledge{
\path{durrett_ch1_composition_measurable};
\path{durrett_ch1_measurability_generator};
\path{durrett_ch1_sigma_field}.
}

\noindent\textbf{Solution.}
\begin{enumerate}[label=(\roman*)]
\item Let $f:\R^d\to\R$ be continuous.  If $G\subseteq\R$ is open, then
$f^{-1}(G)$ is open in $\R^d$, hence belongs to $\B(\R^d)$.  Since the
Borel $\sigma$-field on $\R$ is generated by open sets, $f$ is Borel
measurable.

\item Let $\mathcal{G}$ be a $\sigma$-field on $\R^d$ that makes every
continuous real-valued function measurable.  We show
$\B(\R^d)\subseteq\mathcal{G}$.

For a closed set $C\subseteq\R^d$, define
\[
  d_C(x)=\inf_{y\in C}|x-y|.
\]
The function $d_C$ is continuous, so it is $\mathcal{G}$-measurable.
Since
\[
  C=\{x:d_C(x)=0\}=d_C^{-1}(\{0\}),
\]
we have $C\in\mathcal{G}$.  Thus every closed set belongs to
$\mathcal{G}$, and therefore every open set belongs to $\mathcal{G}$ as
well.  Hence $\B(\R^d)\subseteq\mathcal{G}$.

On the other hand, part (i) shows that $\B(\R^d)$ itself makes all
continuous functions measurable.  Therefore $\B(\R^d)$ is the smallest
such $\sigma$-field.
\end{enumerate}

\section*{Exercise 1.3.5}

\noindent\textbf{Question.}
A function $f$ is lower semicontinuous if
\[
  \liminf_{y\to x} f(y)\ge f(x),
\]
and upper semicontinuous if $-f$ is lower semicontinuous.  Show that
$f$ is lower semicontinuous if and only if $\{x:f(x)\le a\}$ is closed
for each $a\in\R$, and conclude that semicontinuous functions are
measurable.

\Knowledge{
\path{durrett_ch1_measurability_generator};
\path{durrett_ch1_limits_measurable};
\path{durrett_ch1_sigma_field}.
}

\noindent\textbf{Solution.}
Suppose first that $f$ is lower semicontinuous.  Let
$x_n\in\{x:f(x)\le a\}$ and $x_n\to x$.  Then
\[
  f(x)\le \liminf_{n\to\infty} f(x_n)\le a.
\]
Hence $x\in\{f\le a\}$, so $\{f\le a\}$ is closed.

Conversely, assume that $\{f\le a\}$ is closed for every $a\in\R$.
Fix $x$ and take any sequence $y_n\to x$.  If
$\liminf_n f(y_n)<f(x)$, choose $a$ such that
\[
  \liminf_n f(y_n)<a<f(x).
\]
Then infinitely many $y_n$ lie in the closed set $\{f\le a\}$, and a
subsequence of those points converges to $x$.  Closedness gives
$x\in\{f\le a\}$, contradicting $a<f(x)$.  Therefore
\[
  \liminf_{y\to x}f(y)\ge f(x),
\]
so $f$ is lower semicontinuous.

If $f$ is lower semicontinuous, then each set $\{f\le a\}$ is closed and
hence Borel.  This implies measurability because, for example,
\[
  \{f<a\}=\bigcup_{n=1}^\infty \{f\le a-1/n\}
\]
is Borel.  If $f$ is upper semicontinuous, then $-f$ is lower
semicontinuous and measurable, so $f$ is measurable as well.

\section*{Exercise 1.3.6}

\noindent\textbf{Question.}
Let $f:\R^d\to\R$ be arbitrary and define
\[
  f^\delta(x)=\sup\{f(y):|y-x|<\delta\},\qquad
  f_\delta(x)=\inf\{f(y):|y-x|<\delta\}.
\]
Show that $f^\delta$ is lower semicontinuous and $f_\delta$ is upper
semicontinuous.  Let
\[
  f^0=\lim_{\delta\downarrow0}f^\delta,\qquad
  f_0=\lim_{\delta\downarrow0}f_\delta.
\]
Conclude that the set of discontinuities of $f$ is
$\{f^0\ne f_0\}$ and is measurable.

\Knowledge{
\path{durrett_ch1_limits_measurable};
\path{durrett_ch1_measurability_generator};
\path{durrett_ch1_composition_measurable}.
}

\noindent\textbf{Solution.}
We first prove that $f^\delta$ is lower semicontinuous.  It is enough to
show that $\{x:f^\delta(x)>a\}$ is open for each $a\in\R$.  If
$f^\delta(x)>a$, then there is $y$ with $|y-x|<\delta$ and $f(y)>a$.
For all $x'$ sufficiently close to $x$, we still have $|y-x'|<\delta$,
so $f^\delta(x')>a$.  Hence $\{f^\delta>a\}$ is open, and $f^\delta$ is
lower semicontinuous.

Similarly,
\[
  f_\delta(x)=-\sup\{-f(y):|y-x|<\delta\},
\]
so $f_\delta$ is upper semicontinuous.

As $\delta\downarrow0$, the quantities $f^\delta(x)$ decrease and
$f_\delta(x)$ increase, so $f^0$ and $f_0$ are pointwise limits of
measurable functions.  Hence $f^0$ and $f_0$ are measurable.

For every $x$ and every $\delta>0$,
\[
  f_\delta(x)\le f(x)\le f^\delta(x),
\]
so $f_0(x)\le f(x)\le f^0(x)$.  If $f$ is continuous at $x$, then the
oscillation of $f$ on small balls around $x$ tends to $0$, and therefore
$f^0(x)=f_0(x)=f(x)$.  Conversely, if $f^0(x)=f_0(x)$, then the upper
and lower values of $f$ on small balls squeeze to the same number, which
must be $f(x)$; hence $f$ is continuous at $x$.  Therefore the
discontinuity set is exactly
\[
  \{x:f^0(x)\ne f_0(x)\}.
\]
Since $f^0$ and $f_0$ are measurable, this set is measurable.

\section*{Exercise 1.3.7}

\noindent\textbf{Question.}
A function $\phi:\Omega\to\R$ is simple if
\[
  \phi(\omega)=\sum_{m=1}^n c_m\mathbf{1}_{A_m}(\omega),
  \qquad A_m\in\F.
\]
Show that the class of $\F$-measurable functions is the smallest class
containing the simple functions and closed under pointwise limits.

\Knowledge{
\path{durrett_ch1_simple_functions};
\path{durrett_ch1_limits_measurable};
\path{durrett_ch1_measurable_map}.
}

\noindent\textbf{Solution.}
Every simple function is measurable, since for each Borel set $B$,
\[
  \{\phi\in B\}
  =
  \bigcup_{m:c_m\in B} A_m,
\]
a finite union of measurable sets.  Also, pointwise limits of measurable
functions are measurable.  Indeed, if $f_n\to f$, then
\[
  f=\limsup_{n\to\infty} f_n,
\]
and limits, limsups, and liminfs of measurable functions are measurable.
Thus the class of measurable functions contains the simple functions and
is closed under pointwise limits.

It remains to show minimality.  Let $\mathcal{H}$ be any class of
functions that contains the simple functions and is closed under
pointwise limits.  If $f\ge0$ is measurable, define
\[
  f_n(\omega)
  =
  2^{-n}\left\lfloor 2^n\bigl(f(\omega)\wedge n\bigr)\right\rfloor.
\]
Each $f_n$ is a simple measurable function, and $f_n(\omega)\to f(\omega)$
for every $\omega$.  Hence $f\in\mathcal{H}$.

For a general real-valued measurable $f$, write
$f=f^+-f^-$.  Approximate $f^+$ and $f^-$ by simple functions and take
their differences, which are again simple functions.  This gives a
sequence of simple functions converging pointwise to $f$.  Therefore
every measurable function belongs to $\mathcal{H}$.  The measurable
functions are consequently the smallest class containing the simple
functions and closed under pointwise limits.

\section*{Exercise 1.3.8}

\noindent\textbf{Question.}
Use the previous exercise to conclude that $Y$ is measurable with
respect to $\sigma(X)$ if and only if
\[
  Y=f(X)
\]
for some measurable function $f:\R\to\R$.

\Knowledge{
\path{durrett_ch1_sigma_x_factorization};
\path{durrett_ch1_composition_measurable};
\path{durrett_ch1_simple_functions}.
}

\noindent\textbf{Solution.}
If $Y=f(X)$ for a Borel measurable $f$, then $Y$ is
$\sigma(X)$-measurable because
\[
  \{Y\in B\}=\{f(X)\in B\}=\{X\in f^{-1}(B)\}\in\sigma(X).
\]

Conversely, suppose $Y$ is $\sigma(X)$-measurable.  Let $\mathcal{H}$ be
the class of $\sigma(X)$-measurable functions that can be written as
$f(X)$ for some Borel measurable $f$.  This class contains all
$\sigma(X)$-measurable simple functions: if
\[
  Y=\sum_{m=1}^n c_m\mathbf{1}_{A_m}
  \quad\text{and}\quad A_m\in\sigma(X),
\]
then each $A_m=\{X\in B_m\}$ for some Borel set $B_m$, and
\[
  Y=\left(\sum_{m=1}^n c_m\mathbf{1}_{B_m}\right)(X).
\]
The class $\mathcal{H}$ is also closed under pointwise limits, because
if $Y_n=f_n(X)$ and $Y_n\to Y$, then
\[
  Y=f(X),
\]
where $f$ is defined as the pointwise limit of $f_n$ on the Borel set
where that real limit exists, and is defined to be $0$ elsewhere.  This
$f$ is Borel measurable because pointwise limsup and liminf are
measurable.  By Exercise 1.3.7, every $\sigma(X)$-measurable
real-valued function belongs to $\mathcal{H}$.  Thus $Y=f(X)$ for some
Borel measurable $f$.

\section*{Exercise 1.3.9}

\noindent\textbf{Question.}
Give a constructive proof of the previous result.  Note that
\[
  \{\omega:m2^{-n}\le Y(\omega)<(m+1)2^{-n}\}
  =
  \{X\in B_{m,n}\}
\]
for some Borel set $B_{m,n}$, and define $f_n(x)=m2^{-n}$ for
$x\in B_{m,n}$.  Show that $f_n(x)\to f(x)$ and $Y=f(X)$.

\Knowledge{
\path{durrett_ch1_sigma_x_factorization};
\path{durrett_ch1_measurability_generator};
\path{durrett_ch1_limits_measurable}.
}

\noindent\textbf{Solution.}
Assume $Y$ is $\sigma(X)$-measurable.  For each $m\in\mathbb{Z}$ and
$n\ge1$, let
\[
  A_{m,n}
  =
  \{\omega:m2^{-n}\le Y(\omega)<(m+1)2^{-n}\}.
\]
Then $A_{m,n}\in\sigma(X)$, so by the definition of $\sigma(X)$ there is
a Borel set $B_{m,n}\subseteq\R$ such that
\[
  A_{m,n}=\{X\in B_{m,n}\}.
\]
For fixed $n$, the sets $A_{m,n}$ partition $\Omega$.  The corresponding
Borel sets $B_{m,n}$ may overlap away from the actual range of $X$, so
we harmlessly make them disjoint.  Choose an enumeration
$m_1,m_2,\ldots$ of $\mathbb{Z}$ and set
\[
  C_{m_j,n}
  =
  B_{m_j,n}\setminus\bigcup_{i<j}B_{m_i,n}.
\]
Then the $C_{m,n}$ are Borel and disjoint, and still
\[
  A_{m,n}=\{X\in C_{m,n}\}.
\]
Indeed, points removed from $B_{m,n}$ cannot be hit by $X(\omega)$ with
$\omega\in A_{m,n}$, since that would put $\omega$ in two different
partition elements.

Define
\[
  f_n(x)=\sum_{m\in\mathbb{Z}} m2^{-n}\mathbf{1}_{C_{m,n}}(x),
\]
with value $0$ on the complement of $\bigcup_m C_{m,n}$.  Each $f_n$ is
Borel measurable.  Moreover, if $\omega\in A_{m,n}$, then
\[
  f_n(X(\omega))=m2^{-n}
  \quad\text{and}\quad
  m2^{-n}\le Y(\omega)<(m+1)2^{-n}.
\]
Thus
\[
  0\le Y(\omega)-f_n(X(\omega))<2^{-n}
\]
for every $\omega$.  Hence $f_n(X)\to Y$ pointwise.

Now let $D$ be the Borel set on which the real limit
$\lim_{n\to\infty} f_n(x)$ exists, and define
\[
  f(x)=
  \begin{cases}
  \lim_{n\to\infty} f_n(x), & x\in D,\\
  0, & x\notin D.
  \end{cases}
\]
The set $D$ is Borel because it is described by equality of the
measurable functions $\limsup_n f_n$ and $\liminf_n f_n$, and $f$ is
Borel measurable.  Since $f_n(X(\omega))$ converges for every $\omega$,
we have $X(\omega)\in D$.  Evaluating at $X(\omega)$ gives
\[
  f(X(\omega))
  =
  \lim_{n\to\infty} f_n(X(\omega))
  =
  Y(\omega).
\]
Therefore $Y=f(X)$, giving the desired constructive representation.

\section*{Exercise 1.4.1}

\noindent\textbf{Question.}
Show that if $f\ge0$ and
\[
  \int f\,d\mu=0,
\]
then $f=0$ almost everywhere.

\Knowledge{
\path{durrett_ch1_nonnegative_integral};
\path{durrett_ch1_integral_linearity};
\path{durrett_ch1_measure_continuity}.
}

\noindent\textbf{Solution.}
For each $n\ge1$, let
\[
  A_n=\{x:f(x)>1/n\}.
\]
On $A_n$ we have $f\ge (1/n)\mathbf{1}_{A_n}$, so by monotonicity of the
integral,
\[
  0=\int f\,d\mu
  \ge
  \int \frac1n\mathbf{1}_{A_n}\,d\mu
  =
  \frac1n\mu(A_n).
\]
Hence $\mu(A_n)=0$ for every $n$.

Now
\[
  \{x:f(x)>0\}=\bigcup_{n=1}^{\infty} A_n.
\]
This is a countable union of null sets, so it has measure zero.  Thus
$f=0$ almost everywhere.

\section*{Exercise 1.4.2}

\noindent\textbf{Question.}
Let $f\ge0$ and
\[
  E_{n,m}=\left\{x:\frac{m}{2^n}\le f(x)<\frac{m+1}{2^n}\right\}.
\]
Show that, as $n\to\infty$,
\[
  \sum_{m=1}^{\infty}\frac{m}{2^n}\mu(E_{n,m})
  \uparrow
  \int f\,d\mu.
\]

\Knowledge{
\path{durrett_ch1_simple_functions};
\path{durrett_ch1_nonnegative_integral};
\path{durrett_ch1_monotone_convergence}.
}

\noindent\textbf{Solution.}
Define the dyadic lower approximation
\[
  \phi_n(x)=\sum_{m=1}^{\infty}\frac{m}{2^n}\mathbf{1}_{E_{n,m}}(x).
\]
Equivalently,
\[
  \phi_n(x)=2^{-n}\lfloor 2^n f(x)\rfloor
\]
when $f(x)<\infty$, with the natural interpretation for extended
values.  Thus
\[
  0\le \phi_n(x)\le f(x)
  \quad\text{and}\quad
  \phi_n(x)\uparrow f(x)
\]
pointwise as $n\to\infty$.

The integral of $\phi_n$ is exactly
\[
  \int \phi_n\,d\mu
  =
  \sum_{m=1}^{\infty}\frac{m}{2^n}\mu(E_{n,m}),
\]
where the $m=0$ term would contribute $0$ and is therefore omitted.
Since $\phi_n\uparrow f$, the defining monotone approximation property
of the nonnegative integral gives
\[
  \int \phi_n\,d\mu \uparrow \int f\,d\mu.
\]
Hence the displayed sums increase to $\int f\,d\mu$.

\section*{Exercise 1.4.3}

\noindent\textbf{Question.}
Let $g$ be an integrable function on $\R$ and let $\varepsilon>0$.
Show, using the definition of the integral and interval approximation of
measurable sets, that $g$ can be approximated in $L^1$ by a continuous
function.  More explicitly:
\begin{enumerate}[label=(\roman*)]
\item approximate $g$ by a simple function $\phi$ with
$\int |g-\phi|\,dx<\varepsilon$;
\item approximate the sets in $\phi$ by finite unions of intervals to
obtain a step function $q$ with $\int|\phi-q|\,dx<\varepsilon$;
\item round the corners of $q$ to obtain a continuous function $r$ with
$\int|q-r|\,dx<\varepsilon$.
\end{enumerate}

\Knowledge{
\path{durrett_ch1_integrable_function};
\path{durrett_ch1_simple_functions};
\path{durrett_ch1_integral_linearity}.
}

\noindent\textbf{Solution.}
We prove the three approximations in order.

\begin{enumerate}[label=(\roman*)]
\item Since $g$ is integrable, by the construction of the Lebesgue
integral there is a simple function
\[
  \phi=\sum_{k=1}^N b_k\mathbf{1}_{A_k}
\]
such that
\[
  \int_{\R}|g-\phi|\,dx<\varepsilon.
\]
One way to see this is to approximate $g^+$ and $g^-$ from below by
simple functions and then subtract the approximations.

\item By the interval approximation result for Lebesgue measurable sets,
each $A_k$ can be approximated in measure by a finite union $U_k$ of
intervals so well that
\[
  \sum_{k=1}^N |b_k|\,\lambda(A_k\triangle U_k)<\varepsilon,
\]
where $\lambda$ denotes Lebesgue measure.  Put
\[
  \psi=\sum_{k=1}^N b_k\mathbf{1}_{U_k}.
\]
Then
\[
  \int_{\R}|\phi-\psi|\,dx
  \le
  \sum_{k=1}^N |b_k|\,\lambda(A_k\triangle U_k)
  <\varepsilon.
\]
Refining the finitely many interval endpoints, $\psi$ can be rewritten
as a step function
\[
  q=\sum_{j=1}^M c_j\mathbf{1}_{(a_{j-1},a_j)}
\]
with $a_0<a_1<\cdots<a_M$ and
\[
  \int_{\R}|\phi-q|\,dx<\varepsilon.
\]

\item For each interval $(a_{j-1},a_j)$ choose $\delta_j>0$ so small that
$2\delta_j<a_j-a_{j-1}$ and
\[
  \sum_{j=1}^M 2|c_j|\delta_j<\varepsilon.
\]
Define $r_j$ to be continuous, equal to $0$ outside
$(a_{j-1},a_j)$, equal to $c_j$ on
$[a_{j-1}+\delta_j,a_j-\delta_j]$, and linear on the two short boundary
intervals.  Let
\[
  r=\sum_{j=1}^M r_j.
\]
Then $r$ is continuous and differs from $q$ only on boundary intervals
whose total weighted length is controlled by the choice of the
$\delta_j$.  Hence
\[
  \int_{\R}|q-r|\,dx
  \le
  \sum_{j=1}^M 2|c_j|\delta_j
  <\varepsilon.
\]
\end{enumerate}

Combining the three estimates,
\[
  \int_{\R}|g-r|\,dx
  \le
  \int|g-\phi|\,dx+\int|\phi-q|\,dx+\int|q-r|\,dx
  <3\varepsilon.
\]
Since $\varepsilon$ was arbitrary, this proves that integrable functions
can be approximated in $L^1$ by continuous functions.

\section*{Exercise 1.4.4}

\noindent\textbf{Question.}
Prove the Riemann--Lebesgue lemma: if $g$ is integrable, then
\[
  \lim_{n\to\infty}\int_{\R} g(x)\cos(nx)\,dx=0.
\]
Hint: first prove it for step functions, then use the previous exercise.

\Knowledge{
\path{durrett_ch1_integrable_function};
\path{durrett_ch1_integral_linearity};
\path{durrett_ch1_dominated_convergence}.
}

\noindent\textbf{Solution.}
First suppose $q$ is a step function:
\[
  q(x)=\sum_{j=1}^M c_j\mathbf{1}_{(a_{j-1},a_j)}(x).
\]
Then
\[
  \int_{\R}q(x)\cos(nx)\,dx
  =
  \sum_{j=1}^M c_j\int_{a_{j-1}}^{a_j}\cos(nx)\,dx
  =
  \sum_{j=1}^M c_j
  \frac{\sin(na_j)-\sin(na_{j-1})}{n}.
\]
Each term is $O(1/n)$, so this tends to $0$ as $n\to\infty$.

Now let $g\in L^1(\R)$ and fix $\varepsilon>0$.  By Exercise 1.4.3, or
already by its step-function approximation stage, choose a step function
$q$ such that
\[
  \int_{\R}|g-q|\,dx<\varepsilon.
\]
Then
\[
\begin{aligned}
  \left|\int_{\R}g(x)\cos(nx)\,dx\right|
  &\le
  \left|\int_{\R}q(x)\cos(nx)\,dx\right|
  +
  \int_{\R}|g(x)-q(x)|\,|\cos(nx)|\,dx \\
  &\le
  \left|\int_{\R}q(x)\cos(nx)\,dx\right|+\varepsilon.
\end{aligned}
\]
Taking $\limsup_{n\to\infty}$ and using the step-function case gives
\[
  \limsup_{n\to\infty}
  \left|\int_{\R}g(x)\cos(nx)\,dx\right|
  \le \varepsilon.
\]
Since $\varepsilon>0$ is arbitrary, the limit is $0$.

\section*{Exercise 1.5.1}

\noindent\textbf{Question.}
Let
\[
  \lVert f\rVert_\infty
  =
  \inf\{M:\mu(\{x:|f(x)|>M\})=0\}.
\]
Prove that
\[
  \int |fg|\,d\mu\le \lVert f\rVert_1\lVert g\rVert_\infty.
\]

\Knowledge{
\path{durrett_ch1_holder_integral};
\path{durrett_ch1_integrable_function};
\path{durrett_ch1_integral_linearity}.
}

\noindent\textbf{Solution.}
Let $M=\lVert g\rVert_\infty$.  By definition of essential supremum,
$|g|\le M$ almost everywhere.  Therefore
\[
  |fg|\le M|f|
\]
almost everywhere.  By monotonicity and homogeneity of the integral,
\[
  \int |fg|\,d\mu
  \le
  \int M|f|\,d\mu
  =
  M\int |f|\,d\mu
  =
  \lVert g\rVert_\infty\lVert f\rVert_1.
\]
This proves the desired inequality.

\section*{Exercise 1.5.2}

\noindent\textbf{Question.}
Show that if $\mu$ is a probability measure, then
\[
  \lVert f\rVert_\infty=\lim_{p\to\infty}\lVert f\rVert_p.
\]

\Knowledge{
\path{durrett_ch1_jensen_integral};
\path{durrett_ch1_holder_integral};
\path{durrett_ch1_monotone_convergence}.
}

\noindent\textbf{Solution.}
Let $M=\lVert f\rVert_\infty$.  First assume $M<\infty$.  Since
$|f|\le M$ almost everywhere and $\mu(\Omega)=1$,
\[
  \lVert f\rVert_p^p=\int |f|^p\,d\mu\le M^p,
\]
so $\lVert f\rVert_p\le M$ for all $p$.

For the reverse inequality, fix $\varepsilon>0$ and let
\[
  A_\varepsilon=\{|f|>M-\varepsilon\}.
\]
By the definition of essential supremum, $\mu(A_\varepsilon)>0$.  Hence
\[
  \lVert f\rVert_p^p
  \ge
  \int_{A_\varepsilon}|f|^p\,d\mu
  \ge
  (M-\varepsilon)^p\mu(A_\varepsilon).
\]
Taking $p$th roots,
\[
  \lVert f\rVert_p
  \ge
  (M-\varepsilon)\mu(A_\varepsilon)^{1/p}.
\]
Letting $p\to\infty$ gives
\[
  \liminf_{p\to\infty}\lVert f\rVert_p\ge M-\varepsilon.
\]
Since $\varepsilon>0$ is arbitrary and $\limsup_p\lVert f\rVert_p\le M$,
we obtain $\lVert f\rVert_p\to M$.

If $M=\infty$, then for every $K>0$, the set $\{|f|>K\}$ has positive
measure.  The same lower bound gives
\[
  \liminf_{p\to\infty}\lVert f\rVert_p\ge K.
\]
Since $K$ is arbitrary, $\lVert f\rVert_p\to\infty$.  Thus the statement
holds also in the infinite case.

\section*{Exercise 1.5.3}

\noindent\textbf{Question.}
Prove Minkowski's inequality.  For $p\in(1,\infty)$, use Holder's
inequality to show
\[
  \lVert f+g\rVert_p\le \lVert f\rVert_p+\lVert g\rVert_p.
\]
Then show the same result for $p=1$ and $p=\infty$.

\Knowledge{
\path{durrett_ch1_minkowski};
\path{durrett_ch1_holder_integral};
\path{durrett_ch1_integral_linearity}.
}

\noindent\textbf{Solution.}
Let $1<p<\infty$ and let $q$ be the conjugate exponent:
$1/p+1/q=1$.  If $\lVert f+g\rVert_p=0$, there is nothing to prove.
Otherwise,
\[
  \lVert f+g\rVert_p^p
  =
  \int |f+g|^p\,d\mu
  \le
  \int |f|\,|f+g|^{p-1}\,d\mu
  +
  \int |g|\,|f+g|^{p-1}\,d\mu.
\]
By Holder's inequality,
\[
  \int |f|\,|f+g|^{p-1}\,d\mu
  \le
  \lVert f\rVert_p
  \left(\int |f+g|^{(p-1)q}\,d\mu\right)^{1/q}.
\]
Since $(p-1)q=p$, the last factor is
\[
  \left(\int |f+g|^p\,d\mu\right)^{1/q}
  =
  \lVert f+g\rVert_p^{p/q}
  =
  \lVert f+g\rVert_p^{p-1}.
\]
The same estimate holds with $g$ in place of $f$.  Therefore
\[
  \lVert f+g\rVert_p^p
  \le
  \bigl(\lVert f\rVert_p+\lVert g\rVert_p\bigr)
  \lVert f+g\rVert_p^{p-1}.
\]
Dividing by $\lVert f+g\rVert_p^{p-1}$ gives Minkowski's inequality.

For $p=1$,
\[
  \lVert f+g\rVert_1
  =
  \int |f+g|\,d\mu
  \le
  \int |f|\,d\mu+\int |g|\,d\mu
  =
  \lVert f\rVert_1+\lVert g\rVert_1.
\]
For $p=\infty$, $|f+g|\le |f|+|g|$ almost everywhere, hence
\[
  \lVert f+g\rVert_\infty
  \le
  \lVert f\rVert_\infty+\lVert g\rVert_\infty.
\]

\section*{Exercise 1.5.4}

\noindent\textbf{Question.}
If $f$ is integrable and $E_m$ are disjoint sets with union $E$, show
that
\[
  \sum_{m=0}^{\infty}\int_{E_m}f\,d\mu
  =
  \int_E f\,d\mu.
\]
Conclude that if $f\ge0$, then
\[
  \nu(E)=\int_E f\,d\mu
\]
defines a measure.

\Knowledge{
\path{durrett_ch1_integrable_function};
\path{durrett_ch1_integral_linearity};
\path{durrett_ch1_monotone_convergence}.
}

\noindent\textbf{Solution.}
First suppose $f\ge0$.  Since the sets $E_m$ are disjoint,
\[
  \mathbf{1}_E f
  =
  \sum_{m=0}^{\infty}\mathbf{1}_{E_m}f.
\]
Let
\[
  s_n=\sum_{m=0}^n \mathbf{1}_{E_m}f.
\]
Then $s_n\uparrow \mathbf{1}_E f$.  By monotone convergence,
\[
  \int_E f\,d\mu
  =
  \int \mathbf{1}_E f\,d\mu
  =
  \lim_{n\to\infty}\int s_n\,d\mu
  =
  \lim_{n\to\infty}\sum_{m=0}^n\int_{E_m} f\,d\mu.
\]

For a general integrable $f$, write $f=f^+-f^-$.  Applying the
nonnegative case to $f^+$ and $f^-$ and subtracting gives the desired
identity.

If $f\ge0$ and $\nu(E)=\int_E f\,d\mu$, then $\nu(\varnothing)=0$ and
$\nu(E)\ge0$.  The identity above gives countable additivity over
disjoint unions, so $\nu$ is a measure.

\section*{Exercise 1.5.5}

\noindent\textbf{Question.}
If $g_n\uparrow g$ and
\[
  \int g_1^-\,d\mu<\infty,
\]
show that
\[
  \int g_n\,d\mu\uparrow \int g\,d\mu.
\]

\Knowledge{
\path{durrett_ch1_monotone_convergence};
\path{durrett_ch1_integrable_function};
\path{durrett_ch1_integral_linearity}.
}

\noindent\textbf{Solution.}
Since $g_n\uparrow g$, the functions
\[
  h_n=g_n-g_1
\]
are nonnegative and increase pointwise to
\[
  h=g-g_1.
\]
By monotone convergence,
\[
  \int h_n\,d\mu\uparrow \int h\,d\mu.
\]
Since $\int g_1^-\,d\mu<\infty$, the integrals involving $g_1$ are
well-defined in the extended sense, and
\[
  \int g_n\,d\mu
  =
  \int g_1\,d\mu+\int h_n\,d\mu
  \uparrow
  \int g_1\,d\mu+\int h\,d\mu
  =
  \int g\,d\mu.
\]

\section*{Exercise 1.5.6}

\noindent\textbf{Question.}
If $g_m\ge0$, show that
\[
  \int \sum_{m=0}^{\infty}g_m\,d\mu
  =
  \sum_{m=0}^{\infty}\int g_m\,d\mu.
\]

\Knowledge{
\path{durrett_ch1_monotone_convergence};
\path{durrett_ch1_nonnegative_integral};
\path{durrett_ch1_integral_linearity}.
}

\noindent\textbf{Solution.}
Let
\[
  s_n=\sum_{m=0}^n g_m.
\]
Then $s_n\ge0$ and
\[
  s_n\uparrow \sum_{m=0}^{\infty}g_m.
\]
By monotone convergence,
\[
  \int \sum_{m=0}^{\infty}g_m\,d\mu
  =
  \lim_{n\to\infty}\int s_n\,d\mu.
\]
For each finite $n$, linearity gives
\[
  \int s_n\,d\mu
  =
  \sum_{m=0}^n\int g_m\,d\mu.
\]
Taking $n\to\infty$ proves the formula.

\section*{Exercise 1.5.7}

\noindent\textbf{Question.}
Let $f\ge0$.
\begin{enumerate}[label=(\roman*)]
\item Show that
\[
  \int f\wedge n\,d\mu\uparrow \int f\,d\mu.
\]
\item Use this to conclude that if $g$ is integrable and
$\varepsilon>0$, then there is $\delta>0$ such that
$\mu(A)<\delta$ implies
\[
  \int_A |g|\,d\mu<\varepsilon.
\]
\end{enumerate}

\Knowledge{
\path{durrett_ch1_monotone_convergence};
\path{durrett_ch1_integrable_function};
\path{durrett_ch1_integral_linearity}.
}

\noindent\textbf{Solution.}
\begin{enumerate}[label=(\roman*)]
\item Since $f\wedge n\uparrow f$, monotone convergence gives
\[
  \int f\wedge n\,d\mu\uparrow \int f\,d\mu.
\]

\item Apply part (i) to $f=|g|$.  Since $g$ is integrable,
\[
  \int (|g|-|g|\wedge n)\,d\mu\downarrow0.
\]
Choose $n$ such that
\[
  \int (|g|-|g|\wedge n)\,d\mu<\varepsilon/2.
\]
If $A$ is measurable, then
\[
  \int_A |g|\,d\mu
  \le
  \int_A (|g|\wedge n)\,d\mu
  +
  \int (|g|-|g|\wedge n)\,d\mu
  \le
  n\mu(A)+\varepsilon/2.
\]
Now choose $\delta=\varepsilon/(2n)$, with any positive $\delta$ if
$n=0$.  Then $\mu(A)<\delta$ implies
\[
  \int_A |g|\,d\mu<\varepsilon.
\]
\end{enumerate}

\section*{Exercise 1.5.8}

\noindent\textbf{Question.}
Show that if $f$ is integrable on $[a,b]$ and
\[
  g(x)=\int_{[a,x]} f(y)\,dy,
\]
then $g$ is continuous on $(a,b)$.

\Knowledge{
\path{durrett_ch1_integrable_function};
\path{durrett_ch1_integral_linearity};
\path{durrett_ch1_dominated_convergence}.
}

\noindent\textbf{Solution.}
Fix $x\in(a,b)$ and let $\varepsilon>0$.  By Exercise 1.5.7, there is
$\delta>0$ such that whenever a measurable set $A$ has Lebesgue measure
less than $\delta$,
\[
  \int_A |f(y)|\,dy<\varepsilon.
\]
If $h$ is small enough that $x+h\in(a,b)$ and $|h|<\delta$, then
\[
  g(x+h)-g(x)
  =
  \begin{cases}
  \displaystyle\int_{(x,x+h]} f(y)\,dy, & h>0,\\[1ex]
  \displaystyle-\int_{(x+h,x]} f(y)\,dy, & h<0.
  \end{cases}
\]
In either case,
\[
  |g(x+h)-g(x)|
  \le
  \int_{(x\wedge(x+h),\,x\vee(x+h)]}|f(y)|\,dy
  <\varepsilon.
\]
Thus $g(x+h)\to g(x)$ as $h\to0$, so $g$ is continuous on $(a,b)$.

\section*{Exercise 1.5.9}

\noindent\textbf{Question.}
Show that if
\[
  \lVert f\rVert_p=\left(\int |f|^p\,d\mu\right)^{1/p}<\infty,
\]
then there are simple functions $\phi_n$ such that
\[
  \lVert \phi_n-f\rVert_p\to0.
\]

\Knowledge{
\path{durrett_ch1_simple_functions};
\path{durrett_ch1_dominated_convergence};
\path{durrett_ch1_integrable_function}.
}

\noindent\textbf{Solution.}
Assume $1\le p<\infty$.  Since $f$ is measurable, there exists a
sequence of simple functions $\phi_n$ such that
\[
  \phi_n(x)\to f(x)
  \quad\text{and}\quad
  |\phi_n(x)|\le |f(x)|
\]
for all $x$; construct this by applying dyadic simple approximations to
$f^+$ and $f^-$.

Then
\[
  |\phi_n-f|^p
  \le
  (|\phi_n|+|f|)^p
  \le
  (2|f|)^p
  =
  2^p|f|^p.
\]
The function $2^p|f|^p$ is integrable because $\lVert f\rVert_p<\infty$.
Also $|\phi_n-f|^p\to0$ pointwise.  By dominated convergence,
\[
  \int |\phi_n-f|^p\,d\mu\to0.
\]
Therefore $\lVert \phi_n-f\rVert_p\to0$.

\section*{Exercise 1.5.10}

\noindent\textbf{Question.}
Show that if
\[
  \sum_n \int |f_n|\,d\mu<\infty,
\]
then
\[
  \sum_n\int f_n\,d\mu
  =
  \int \sum_n f_n\,d\mu.
\]

\Knowledge{
\path{durrett_ch1_dominated_convergence};
\path{durrett_ch1_monotone_convergence};
\path{durrett_ch1_integral_linearity}.
}

\noindent\textbf{Solution.}
By Exercise 1.5.6 applied to the nonnegative functions $|f_n|$,
\[
  \int \sum_n |f_n|\,d\mu
  =
  \sum_n\int |f_n|\,d\mu
  <\infty.
\]
Hence
\[
  H=\sum_n |f_n|
\]
is integrable and finite almost everywhere.  Therefore
\[
  \sum_n f_n
\]
converges absolutely almost everywhere.

Let
\[
  S_N=\sum_{n=0}^N f_n.
\]
Then $S_N\to S=\sum_n f_n$ almost everywhere and
\[
  |S_N|\le H
\]
for all $N$.  By dominated convergence,
\[
  \int S_N\,d\mu\to \int S\,d\mu.
\]
For finite sums, linearity gives
\[
  \int S_N\,d\mu=\sum_{n=0}^N\int f_n\,d\mu.
\]
Letting $N\to\infty$ yields
\[
  \int \sum_n f_n\,d\mu
  =
  \sum_n\int f_n\,d\mu.
\]

\section*{Exercise 1.6.1}

\noindent\textbf{Question.}
Suppose $\phi$ is strictly convex.  Under the assumptions of Jensen's
inequality, show that
\[
  \phi(EX)=E\phi(X)
\]
implies $X=EX$ almost surely.

\Knowledge{
\path{durrett_ch1_jensen_integral};
\path{durrett_ch1_expectation_as_integral}.
}

\noindent\textbf{Solution.}
Let $m=EX$.  If $P(X=m)=1$, equality is clear.  Conversely suppose
$P(X\ne m)>0$.  Then there are numbers $a<m<b$ such that
\[
  P(X\le a)>0,\qquad P(X\ge b)>0.
\]
By strict convexity, the graph of $\phi$ lies strictly above any
supporting chord except at the endpoints.  Equivalently, Jensen's
inequality is strict for a nonconstant random variable.  More explicitly,
condition on the two sets $\{X<m\}$ and $\{X\ge m\}$ and apply strict
convexity to the two conditional means; the inequality is strict because
both sets have positive probability and their conditional means are
different.  Thus
\[
  E\phi(X)>\phi(EX),
\]
contradicting equality.  Hence $P(X=EX)=1$.

\section*{Exercise 1.6.2}

\noindent\textbf{Question.}
Suppose $\phi:\R^n\to\R$ is convex.  Show that
\[
  E\phi(X_1,\ldots,X_n)\ge \phi(EX_1,\ldots,EX_n),
\]
provided $E|\phi(X_1,\ldots,X_n)|<\infty$ and $E|X_i|<\infty$ for all
$i$.

\Knowledge{
\path{durrett_ch1_jensen_integral};
\path{durrett_ch1_expectation_linearity}.
}

\noindent\textbf{Solution.}
Let
\[
  m=(EX_1,\ldots,EX_n).
\]
For a convex function on $\R^n$, there is a supporting hyperplane at
$m$: there are $a\in\R^n$ and $b\in\R$ such that
\[
  \phi(x)\ge a\cdot x+b\quad\text{for all }x,
  \qquad
  \phi(m)=a\cdot m+b.
\]
Therefore
\[
  E\phi(X_1,\ldots,X_n)
  \ge
  E\bigl[a\cdot(X_1,\ldots,X_n)+b\bigr]
  =
  a\cdot m+b
  =
  \phi(m).
\]
This proves the multidimensional Jensen inequality.

\section*{Exercise 1.6.3}

\noindent\textbf{Question.}
Chebyshev's inequality is and is not sharp.
\begin{enumerate}[label=(\roman*)]
\item Show that if $0<b\le a$ are fixed, there is $X$ with $EX^2=b^2$
for which $P(|X|\ge a)=b^2/a^2$.
\item Show that if $0<EX^2<\infty$, then
\[
  \lim_{a\to\infty}\frac{a^2P(|X|\ge a)}{EX^2}=0.
\]
\end{enumerate}

\Knowledge{
\path{durrett_ch1_markov_chebyshev};
\path{durrett_ch1_expectation_limit_theorems}.
}

\noindent\textbf{Solution.}
\begin{enumerate}[label=(\roman*)]
\item Let
\[
  P(X=a)=\frac{b^2}{2a^2},\qquad
  P(X=-a)=\frac{b^2}{2a^2},\qquad
  P(X=0)=1-\frac{b^2}{a^2}.
\]
Then
\[
  EX^2=a^2\frac{b^2}{a^2}=b^2
  \quad\text{and}\quad
  P(|X|\ge a)=\frac{b^2}{a^2}.
\]
So Chebyshev's bound can be attained at a fixed level $a$.

\item Let $Y=X^2$.  Since $EY<\infty$,
\[
  a^2P(|X|\ge a)=a^2P(Y\ge a^2)
  \le E(Y;Y\ge a^2).
\]
As $a\to\infty$, $Y\mathbf{1}_{\{Y\ge a^2\}}\downarrow0$ and is bounded
by the integrable random variable $Y$.  By dominated convergence,
\[
  E(Y;Y\ge a^2)\to0.
\]
Dividing by the positive constant $EX^2$ gives the desired limit.
\end{enumerate}

\section*{Exercise 1.6.4}

\noindent\textbf{Question.}
One-sided Chebyshev bound.
\begin{enumerate}[label=(\roman*)]
\item Let $a>b>0$, $0<p<1$, and let $X$ satisfy
$P(X=a)=p$, $P(X=-b)=1-p$.  Apply Chebyshev's inequality to
$\phi(x)=(x+b)^2$ and conclude that if $Y$ has the same mean and
variance as $X$, then $P(Y\ge a)\le p$, with equality when $Y=X$.
\item If $EY=0$, $\operatorname{var}(Y)=\sigma^2$, and $a>0$, show
\[
  P(Y\ge a)\le \frac{\sigma^2}{a^2+\sigma^2},
\]
and show equality is possible.
\end{enumerate}

\Knowledge{
\path{durrett_ch1_markov_chebyshev};
\path{durrett_ch1_expectation_linearity};
\path{durrett_ch1_common_distributions_expectations}.
}

\noindent\textbf{Solution.}
\begin{enumerate}[label=(\roman*)]
\item If $X$ has the stated distribution, then
\[
  EX=pa-(1-p)b,\qquad E(X+b)^2=p(a+b)^2.
\]
For any $Y$ with the same mean and variance as $X$, we also have
$E(Y+b)^2=E(X+b)^2=p(a+b)^2$, because $E(Y+b)^2$ is determined by
$EY$ and $EY^2$.  On the event $\{Y\ge a\}$,
\[
  (Y+b)^2\ge(a+b)^2.
\]
Thus Chebyshev's inequality gives
\[
  (a+b)^2P(Y\ge a)\le E(Y+b)^2=p(a+b)^2,
\]
so $P(Y\ge a)\le p$.  For $Y=X$, equality holds.

\item Apply part (i) with a two-point random variable taking values
$a$ and $-b$ and having mean $0$.  The mean-zero condition gives
\[
  pa-(1-p)b=0,\qquad p=\frac{b}{a+b}.
\]
Its variance is
\[
  \sigma^2=pa^2+(1-p)b^2=ab.
\]
Thus $b=\sigma^2/a$ and
\[
  p=\frac{\sigma^2/a}{a+\sigma^2/a}
  =
  \frac{\sigma^2}{a^2+\sigma^2}.
\]
The first part gives
\[
  P(Y\ge a)\le \frac{\sigma^2}{a^2+\sigma^2}.
\]
Equality is attained by the two-point variable with
$P(Y=a)=\sigma^2/(a^2+\sigma^2)$ and
$P(Y=-\sigma^2/a)=a^2/(a^2+\sigma^2)$.
\end{enumerate}

\section*{Exercise 1.6.5}

\noindent\textbf{Question.}
Show that the following lower bounds do not exist:
\begin{enumerate}[label=(\roman*)]
\item If $\varepsilon>0$,
\[
  \inf\{P(|X|>\varepsilon):EX=0,\operatorname{var}(X)=1\}=0.
\]
\item If $y\ge1$ and $\sigma^2\in(0,\infty)$,
\[
  \inf\{P(|X|>y):EX=1,\operatorname{var}(X)=\sigma^2\}=0.
\]
\end{enumerate}

\Knowledge{
\path{durrett_ch1_common_distributions_expectations};
\path{durrett_ch1_markov_chebyshev}.
}

\noindent\textbf{Solution.}
\begin{enumerate}[label=(\roman*)]
\item Fix $\delta>0$.  Let
\[
  P(X=0)=1-\delta,\qquad
  P(X=\delta^{-1/2})=P(X=-\delta^{-1/2})=\delta/2.
\]
Then $EX=0$ and $EX^2=1$, so $\operatorname{var}(X)=1$.  Also
$P(|X|>\varepsilon)\le\delta$ for all sufficiently small $\delta$ such
that $\delta^{-1/2}>\varepsilon$.  Since $\delta$ is arbitrary, the
infimum is $0$.

\item Fix $\delta>0$ and let $Z$ have mean $0$ and variance $1$ with
$P(Z=0)=1-\delta$ as in part (i).  Put
\[
  X=1+\sigma Z.
\]
Then $EX=1$ and $\operatorname{var}(X)=\sigma^2$.  The event
$\{|X|>y\}$ is contained in the event $\{Z\ne0\}$ unless one of the rare
nonzero values occurs; hence
\[
  P(|X|>y)\le P(Z\ne0)=\delta.
\]
Letting $\delta\downarrow0$ gives the infimum $0$.
\end{enumerate}

\section*{Exercise 1.6.6}

\noindent\textbf{Question.}
Let $Y\ge0$ with $EY^2<\infty$.  Apply Cauchy--Schwarz to
$Y\mathbf{1}_{\{Y>0\}}$ and conclude
\[
  P(Y>0)\ge \frac{(EY)^2}{EY^2}.
\]

\Knowledge{
\path{durrett_ch1_holder_integral};
\path{durrett_ch1_expectation_linearity}.
}

\noindent\textbf{Solution.}
Since $Y=Y\mathbf{1}_{\{Y>0\}}$,
\[
  EY=E\bigl(Y\mathbf{1}_{\{Y>0\}}\bigr).
\]
By Cauchy--Schwarz,
\[
  (EY)^2
  \le
  E(Y^2)\,E(\mathbf{1}_{\{Y>0\}}^2)
  =
  EY^2\,P(Y>0).
\]
If $EY^2>0$, division gives
\[
  P(Y>0)\ge \frac{(EY)^2}{EY^2}.
\]
If $EY^2=0$, then $Y=0$ a.s. and the assertion is trivial.

\section*{Exercise 1.6.7}

\noindent\textbf{Question.}
Let $\Omega=(0,1)$ with Lebesgue measure, let $\alpha\in(1,2)$, and set
\[
  X_n=n^\alpha\mathbf{1}_{(1/(n+1),1/n)}.
\]
Show $X_n\to0$ almost surely and Theorem 1.6.8 applies with
$h(x)=x$ and $g(x)=|x|^{2/\alpha}$, but the $X_n$ are not dominated by
an integrable random variable.

\Knowledge{
\path{durrett_ch1_expectation_limit_theorems};
\path{durrett_ch1_dominated_convergence};
\path{durrett_ch1_common_distributions_expectations}.
}

\noindent\textbf{Solution.}
The intervals $(1/(n+1),1/n)$ are disjoint.  Hence for each
$\omega\in(0,1)$, at most one $X_n(\omega)$ is nonzero.  Thus
$X_n(\omega)\to0$ for every $\omega$.

For $g(x)=|x|^{2/\alpha}$,
\[
  Eg(X_n)
  =
  n^2\,\lambda(1/(n+1),1/n)
  =
  n^2\left(\frac1n-\frac1{n+1}\right)
  =
  \frac{n}{n+1}
  \le1.
\]
Also $g(x)\to\infty$ as $|x|\to\infty$, and
\[
  \frac{|h(x)|}{g(x)}
  =
  |x|^{1-2/\alpha}\to0
  \quad\text{as } |x|\to\infty,
\]
because $\alpha<2$.  Thus Theorem 1.6.8 applies.

Suppose, toward a contradiction, that $|X_n|\le Y$ for all $n$ with
$EY<\infty$.  Then on $(1/(n+1),1/n)$, $Y\ge n^\alpha$.  Therefore
\[
  EY\ge
  \sum_{n=1}^{N} n^\alpha
  \left(\frac1n-\frac1{n+1}\right)
  =
  \sum_{n=1}^{N}\frac{n^{\alpha-1}}{n+1}.
\]
The terms are comparable to $n^{\alpha-2}$, and since
$\alpha-2\in(-1,0)$, the series diverges.  This contradicts
$EY<\infty$.  Hence no integrable dominating random variable exists.

\section*{Exercise 1.6.8}

\noindent\textbf{Question.}
Suppose a probability measure $\mu$ has
\[
  \mu(A)=\int_A f(x)\,dx
\]
for all Borel $A$.  Show that for any $g$ with $g\ge0$ or
$\int |g|\,d\mu<\infty$,
\[
  \int g(x)\,\mu(dx)=\int g(x)f(x)\,dx.
\]

\Knowledge{
\path{durrett_ch1_change_of_variables};
\path{durrett_ch1_expectation_limit_theorems}.
}

\noindent\textbf{Solution.}
First take $g=\mathbf{1}_A$.  Then
\[
  \int \mathbf{1}_A\,d\mu=\mu(A)=\int_A f(x)\,dx
  =
  \int \mathbf{1}_A(x)f(x)\,dx.
\]
By linearity, the result holds for nonnegative simple functions.  If
$g\ge0$, choose simple $g_n\uparrow g$.  Monotone convergence on both
sides gives
\[
  \int g\,d\mu
  =
  \lim_n\int g_n\,d\mu
  =
  \lim_n\int g_n f\,dx
  =
  \int gf\,dx.
\]
For integrable signed $g$, apply the nonnegative result to $g^+$ and
$g^-$ and subtract.

\section*{Exercise 1.6.9}

\noindent\textbf{Question.}
Let $A_1,\ldots,A_n$ be events and $A=\bigcup_{i=1}^n A_i$.  Prove the
inclusion--exclusion formula.

\Knowledge{
\path{durrett_ch1_inclusion_exclusion};
\path{durrett_ch1_expectation_linearity}.
}

\noindent\textbf{Solution.}
For each $\omega$,
\[
  \mathbf{1}_A(\omega)
  =
  1-\prod_{i=1}^n(1-\mathbf{1}_{A_i}(\omega)).
\]
Indeed, the product is $1$ exactly when none of the $A_i$ occur, and is
$0$ otherwise.  Expanding the product gives
\[
  \mathbf{1}_A
  =
  \sum_i \mathbf{1}_{A_i}
  -
  \sum_{i<j}\mathbf{1}_{A_i\cap A_j}
  +
  \sum_{i<j<k}\mathbf{1}_{A_i\cap A_j\cap A_k}
  -\cdots
  +(-1)^{n-1}\mathbf{1}_{\cap_{i=1}^n A_i}.
\]
Taking expectations yields
\[
  P\left(\bigcup_{i=1}^n A_i\right)
  =
  \sum_i P(A_i)
  -
  \sum_{i<j}P(A_i\cap A_j)
  +
  \sum_{i<j<k}P(A_i\cap A_j\cap A_k)
  -\cdots
  +(-1)^{n-1}P\left(\bigcap_{i=1}^n A_i\right).
\]

\section*{Exercise 1.6.10}

\noindent\textbf{Question.}
Prove the Bonferroni inequalities for
$A=\bigcup_{i=1}^n A_i$: stopping inclusion--exclusion after an even
number of sums gives a lower bound, and stopping after an odd number of
sums gives an upper bound.

\Knowledge{
\path{durrett_ch1_inclusion_exclusion};
\path{durrett_ch1_expectation_linearity}.
}

\noindent\textbf{Solution.}
For a fixed outcome $\omega$, let $N(\omega)=\sum_i\mathbf{1}_{A_i}(\omega)$.
Then $\mathbf{1}_A(\omega)=\mathbf{1}_{\{N(\omega)\ge1\}}$.  The
inclusion--exclusion partial sums are the binomial sums
\[
  S_r(N)=\sum_{k=1}^{r}(-1)^{k-1}\binom{N}{k}.
\]
For $N=0$, all sums are $0$.  For $N\ge1$, the full alternating sum is
$1$, and the alternating partial sums lie alternately above and below
$1$:
\[
  S_1(N)\ge1,\qquad S_2(N)\le1,\qquad S_3(N)\ge1,\quad\ldots.
\]
Thus, pointwise,
\[
  \mathbf{1}_A
  \le \sum_i\mathbf{1}_{A_i},
\]
\[
  \mathbf{1}_A
  \ge
  \sum_i\mathbf{1}_{A_i}
  -
  \sum_{i<j}\mathbf{1}_{A_i\cap A_j},
\]
and
\[
  \mathbf{1}_A
  \le
  \sum_i\mathbf{1}_{A_i}
  -
  \sum_{i<j}\mathbf{1}_{A_i\cap A_j}
  +
  \sum_{i<j<k}\mathbf{1}_{A_i\cap A_j\cap A_k}.
\]
Taking expectations gives the displayed Bonferroni inequalities.  The
same pointwise alternating-sum argument gives the general even/lower and
odd/upper rule.

\section*{Exercise 1.6.11}

\noindent\textbf{Question.}
If $E|X|^k<\infty$, show that for $0<j<k$,
\[
  E|X|^j<\infty,
  \qquad
  E|X|^j\le (E|X|^k)^{j/k}.
\]

\Knowledge{
\path{durrett_ch1_jensen_integral};
\path{durrett_ch1_holder_integral}.
}

\noindent\textbf{Solution.}
Let $Y=|X|^k\ge0$ and let $r=j/k\in(0,1)$.  The function
$t\mapsto t^r$ is concave on $[0,\infty)$.  Jensen's inequality for the
convex function $-t^r$ gives
\[
  EY^r\le (EY)^r.
\]
Since $Y^r=|X|^j$, we get
\[
  E|X|^j\le (E|X|^k)^{j/k}<\infty.
\]

\section*{Exercise 1.6.12}

\noindent\textbf{Question.}
Apply Jensen's inequality with $\phi(x)=e^x$ and
$P(X=\log y_m)=p(m)$ to prove
\[
  \sum_{m=1}^n p(m)y_m
  \ge
  \prod_{m=1}^n y_m^{p(m)}
\]
when $\sum_m p(m)=1$ and $p(m),y_m>0$.

\Knowledge{
\path{durrett_ch1_jensen_integral};
\path{durrett_ch1_common_distributions_expectations}.
}

\noindent\textbf{Solution.}
Let $X$ be the discrete random variable with
$P(X=\log y_m)=p(m)$.  Since $\phi(x)=e^x$ is convex, Jensen's
inequality gives
\[
  E(e^X)\ge e^{EX}.
\]
Now
\[
  E(e^X)=\sum_{m=1}^n p(m)y_m,
  \qquad
  e^{EX}
  =
  \exp\left(\sum_{m=1}^n p(m)\log y_m\right)
  =
  \prod_{m=1}^n y_m^{p(m)}.
\]
This proves the weighted arithmetic-geometric mean inequality.  When
$p(m)=1/n$, it is the usual AM-GM inequality.

\section*{Exercise 1.6.13}

\noindent\textbf{Question.}
If $EX_1^-<\infty$ and $X_n\uparrow X$, show that
\[
  EX_n\uparrow EX.
\]

\Knowledge{
\path{durrett_ch1_expectation_limit_theorems};
\path{durrett_ch1_expectation_linearity}.
}

\noindent\textbf{Solution.}
Let
\[
  Y_n=X_n-X_1.
\]
Then $Y_n\ge0$ and $Y_n\uparrow X-X_1$.  By monotone convergence,
\[
  EY_n\uparrow E(X-X_1).
\]
Since $EX_1^-<\infty$, adding $EX_1$ is legitimate in the extended
sense.  Therefore
\[
  EX_n=EX_1+EY_n\uparrow EX_1+E(X-X_1)=EX.
\]

\section*{Exercise 1.6.14}

\noindent\textbf{Question.}
Let $X\ge0$, but do not assume $E(1/X)<\infty$.  Show
\[
  \lim_{y\to\infty} yE(1/X;X>y)=0,
  \qquad
  \lim_{y\downarrow0} yE(1/X;X>y)=0.
\]

\Knowledge{
\path{durrett_ch1_tail_sum_nonnegative};
\path{durrett_ch1_expectation_limit_theorems}.
}

\noindent\textbf{Solution.}
For $y\to\infty$, on $\{X>y\}$ we have $y/X\le1$, so
\[
  0\le yE(1/X;X>y)=E(y/X;X>y)\le P(X>y)\to0.
\]

For $y\downarrow0$, fix $\eta>0$.  Choose $\delta>0$ so small that
\[
  P(0<X\le \delta)<\eta,
\]
which is possible by continuity of probability from above.  If
$0<y<\delta$, then
\[
\begin{aligned}
  yE(1/X;X>y)
  &=
  E(y/X;y<X\le\delta)+E(y/X;X>\delta)\\
  &\le
  P(y<X\le\delta)+\frac{y}{\delta}
  \le
  \eta+\frac{y}{\delta}.
\end{aligned}
\]
Letting $y\downarrow0$ gives a limsup at most $\eta$.  Since $\eta$ is
arbitrary, the second limit is $0$.

\section*{Exercise 1.6.15}

\noindent\textbf{Question.}
If $X_n\ge0$, show that
\[
  E\left(\sum_{n=0}^{\infty}X_n\right)=\sum_{n=0}^{\infty}EX_n.
\]

\Knowledge{
\path{durrett_ch1_expectation_limit_theorems};
\path{durrett_ch1_fubini_tonelli}.
}

\noindent\textbf{Solution.}
Let
\[
  S_N=\sum_{n=0}^N X_n.
\]
Then $S_N\uparrow \sum_{n=0}^{\infty}X_n$.  By monotone convergence,
\[
  E\left(\sum_{n=0}^{\infty}X_n\right)
  =
  \lim_{N\to\infty}ES_N.
\]
For finite sums, expectation is linear, so
\[
  ES_N=\sum_{n=0}^N EX_n.
\]
Letting $N\to\infty$ proves the identity.

\section*{Exercise 1.6.16}

\noindent\textbf{Question.}
If $X$ is integrable and $A_n$ are disjoint sets with union $A$, show
\[
  \sum_{n=0}^{\infty}E(X;A_n)=E(X;A),
\]
i.e. the sum converges absolutely and has the displayed value.

\Knowledge{
\path{durrett_ch1_expectation_as_integral};
\path{durrett_ch1_expectation_linearity};
\path{durrett_ch1_expectation_limit_theorems}.
}

\noindent\textbf{Solution.}
Since $X$ is integrable,
\[
  \sum_{n=0}^{\infty}E(|X|;A_n)
  =
  E\left(|X|;\bigcup_{n=0}^{\infty}A_n\right)
  =
  E(|X|;A)
  \le E|X|<\infty,
\]
where we used Exercise 1.6.15 for the nonnegative random variables
$|X|\mathbf{1}_{A_n}$.  Hence the series
$\sum_n E(X;A_n)$ converges absolutely.

Now use finite additivity and pass to the limit:
\[
  \sum_{n=0}^{N}E(X;A_n)
  =
  E\left(X;\bigcup_{n=0}^{N}A_n\right).
\]
As $N\to\infty$,
\[
  X\mathbf{1}_{\cup_{n=0}^{N}A_n}\to X\mathbf{1}_A
  \quad\text{and}\quad
  \left|X\mathbf{1}_{\cup_{n=0}^{N}A_n}\right|\le |X|.
\]
Dominated convergence gives
\[
  E\left(X;\bigcup_{n=0}^{N}A_n\right)\to E(X;A).
\]
Therefore
\[
  \sum_{n=0}^{\infty}E(X;A_n)=E(X;A).
\]

\section*{Exercise 1.7.1}

\noindent\textbf{Question.}
Show that if
\[
  \int_X\int_Y |f(x,y)|\,\mu_2(dy)\,\mu_1(dx)<\infty,
\]
then
\[
  \int_X\int_Y f(x,y)\,\mu_2(dy)\,\mu_1(dx)
  =
  \int_{X\times Y} f\,d(\mu_1\times\mu_2)
  =
  \int_Y\int_X f(x,y)\,\mu_1(dx)\,\mu_2(dy).
\]

\Knowledge{
\path{durrett_ch1_product_measure};
\path{durrett_ch1_fubini_tonelli};
\path{durrett_ch1_integrable_function}.
}

\noindent\textbf{Solution.}
The hypothesis says that one iterated integral of $|f|$ is finite.  By
Tonelli's theorem applied to the nonnegative function $|f|$,
\[
  \int_{X\times Y}|f|\,d(\mu_1\times\mu_2)
  =
  \int_X\int_Y |f(x,y)|\,\mu_2(dy)\,\mu_1(dx)
  <\infty.
\]
Thus $f$ is integrable with respect to the product measure.  Fubini's
theorem applies to $f$, and gives both iterated integrals and their
equality to the product integral:
\[
  \int_X\int_Y f\,d\mu_2\,d\mu_1
  =
  \int_{X\times Y} f\,d(\mu_1\times\mu_2)
  =
  \int_Y\int_X f\,d\mu_1\,d\mu_2.
\]

The stated corollary follows by taking $X=\{1,2,\ldots\}$ with counting
measure and considering $f(n,\omega)=f_n(\omega)$.  The assumption
$\sum_n\int |f_n|\,d\mu<\infty$ is exactly the integrability hypothesis
above, so
\[
  \sum_n\int f_n\,d\mu=\int\sum_n f_n\,d\mu.
\]

\section*{Exercise 1.7.2}

\noindent\textbf{Question.}
Let $g\ge0$ be measurable on $(X,\mathcal{A},\mu)$.  Use Fubini/Tonelli
to show
\[
  \int_X g\,d\mu
  =
  (\mu\times\lambda)(\{(x,y):0\le y<g(x)\})
  =
  \int_0^\infty \mu(\{x:g(x)>y\})\,dy.
\]

\Knowledge{
\path{durrett_ch1_fubini_tonelli};
\path{durrett_ch1_sections_measurable};
\path{durrett_ch1_tail_sum_nonnegative}.
}

\noindent\textbf{Solution.}
Let
\[
  E=\{(x,y):0\le y<g(x)\}.
\]
Since $g$ is measurable, $E$ is measurable in the product
$\sigma$-field.  For each fixed $x$,
\[
  \lambda(E_x)=\lambda(\{y:0\le y<g(x)\})=g(x).
\]
Therefore Tonelli's theorem gives
\[
  (\mu\times\lambda)(E)
  =
  \int_X \lambda(E_x)\,\mu(dx)
  =
  \int_X g(x)\,\mu(dx).
\]
For each fixed $y\ge0$, the $y$-section is
\[
  E^y=\{x:g(x)>y\}.
\]
Thus another application of Tonelli gives
\[
  (\mu\times\lambda)(E)
  =
  \int_0^\infty \mu(\{x:g(x)>y\})\,dy.
\]
Combining the two identities proves the formula.

\section*{Exercise 1.7.3}

\noindent\textbf{Question.}
Let $F,G$ be Stieltjes measure functions and let $\mu,\nu$ be the
corresponding measures.  Show:
\begin{enumerate}[label=(\roman*)]
\item
\[
  \int_{(a,b]}(F(y)-F(a))\,dG(y)
  =
  (\mu\times\nu)(\{(x,y):a<x\le y\le b\}).
\]
\item
\[
  \int_{(a,b]}F(y)\,dG(y)+\int_{(a,b]}G(y)\,dF(y)
  =
  F(b)G(b)-F(a)G(a)
  +
  \sum_{x\in(a,b]}\mu(\{x\})\nu(\{x\}).
\]
\item If $F=G$ is continuous, then
\[
  \int_{(a,b]}2F(y)\,dF(y)=F^2(b)-F^2(a).
\]
\end{enumerate}

\Knowledge{
\path{durrett_ch1_stieltjes_measure};
\path{durrett_ch1_product_measure};
\path{durrett_ch1_fubini_tonelli}.
}

\noindent\textbf{Solution.}
\begin{enumerate}[label=(\roman*)]
\item For $y\in(a,b]$,
\[
  F(y)-F(a)=\mu((a,y]).
\]
Hence, by Tonelli,
\[
\begin{aligned}
  \int_{(a,b]}(F(y)-F(a))\,dG(y)
  &=
  \int_{(a,b]}\int_{\R}\mathbf{1}_{\{a<x\le y\}}\,\mu(dx)\,\nu(dy)\\
  &=
  (\mu\times\nu)(\{(x,y):a<x\le y\le b\}).
\end{aligned}
\]

\item Let
\[
  D=\{(x,y):a<x\le y\le b\}.
\]
Part (i) gives
\[
  (\mu\times\nu)(D)
  =
  \int_{(a,b]}F(y)\,dG(y)-F(a)(G(b)-G(a)).
\]
Similarly,
\[
  (\mu\times\nu)(D')
  =
  \int_{(a,b]}G(x)\,dF(x)-G(a)(F(b)-F(a)),
\]
where
\[
  D'=\{(x,y):a<y\le x\le b\}.
\]
The two sets $D$ and $D'$ cover the rectangle $(a,b]\times(a,b]$ and
overlap exactly on the diagonal $\{(x,x):x\in(a,b]\}$.  Therefore
\[
  (\mu\times\nu)(D)+(\mu\times\nu)(D')
  =
  (F(b)-F(a))(G(b)-G(a))
  +
  \sum_{x\in(a,b]}\mu(\{x\})\nu(\{x\}).
\]
Substituting the two integral expressions and simplifying gives the
claimed identity.

\item If $F=G$ is continuous, then $\mu(\{x\})=0$ for every $x$, so the
jump sum in part (ii) is zero.  Setting $G=F$ gives
\[
  2\int_{(a,b]}F(y)\,dF(y)=F^2(b)-F^2(a).
\]
\end{enumerate}

\section*{Exercise 1.7.4}

\noindent\textbf{Question.}
Let $\mu$ be a finite measure on $\R$ and
\[
  F(x)=\mu((-\infty,x]).
\]
Show that, for $c>0$,
\[
  \int_{\R}(F(x+c)-F(x))\,dx=c\,\mu(\R).
\]

\Knowledge{
\path{durrett_ch1_stieltjes_measure};
\path{durrett_ch1_fubini_tonelli};
\path{durrett_ch1_integrable_function}.
}

\noindent\textbf{Solution.}
For each $x$,
\[
  F(x+c)-F(x)=\mu((x,x+c]).
\]
Therefore
\[
  F(x+c)-F(x)
  =
  \int_{\R}\mathbf{1}_{(x,x+c]}(t)\,\mu(dt).
\]
By Tonelli's theorem,
\[
\begin{aligned}
  \int_{\R}(F(x+c)-F(x))\,dx
  &=
  \int_{\R}\int_{\R}\mathbf{1}_{(x,x+c]}(t)\,\mu(dt)\,dx\\
  &=
  \int_{\R}\lambda(\{x:x<t\le x+c\})\,\mu(dt).
\end{aligned}
\]
For fixed $t$, the set $\{x:x<t\le x+c\}$ is $[t-c,t)$, which has
Lebesgue measure $c$.  Hence
\[
  \int_{\R}(F(x+c)-F(x))\,dx
  =
  \int_{\R}c\,\mu(dt)
  =
  c\,\mu(\R).
\]

\section*{Exercise 1.7.5}

\noindent\textbf{Question.}
Show that $e^{-xy}\sin x$ is integrable on the strip
$0<x<a$, $0<y<\infty$.  Evaluate the double integral in both orders and
obtain the usual sine-integral estimate.

\Knowledge{
\path{durrett_ch1_fubini_tonelli};
\path{durrett_ch1_integrable_function};
\path{durrett_ch1_dominated_convergence}.
}

\noindent\textbf{Solution.}
First,
\[
  \int_0^a\int_0^\infty e^{-xy}|\sin x|\,dy\,dx
  =
  \int_0^a \frac{|\sin x|}{x}\,dx<\infty,
\]
because $|\sin x|/x$ is bounded near $0$ and integrable on $[0,a]$.
Thus Fubini's theorem applies.

Integrating first in $y$ gives
\[
  \int_0^a\int_0^\infty e^{-xy}\sin x\,dy\,dx
  =
  \int_0^a \frac{\sin x}{x}\,dx.
\]
Integrating first in $x$, we use
\[
  \int_0^a e^{-xy}\sin x\,dx
  =
  \frac{1-e^{-ay}(y\sin a+\cos a)}{1+y^2}.
\]
Therefore
\[
  \int_0^a \frac{\sin x}{x}\,dx
  =
  \int_0^\infty \frac{dy}{1+y^2}
  -
  (\cos a)\int_0^\infty\frac{e^{-ay}}{1+y^2}\,dy
  -
  (\sin a)\int_0^\infty\frac{ye^{-ay}}{1+y^2}\,dy.
\]
Since $\int_0^\infty(1+y^2)^{-1}\,dy=\pi/2$, this is
\[
  \int_0^a \frac{\sin x}{x}\,dx
  =
  \frac{\pi}{2}
  -
  (\cos a)\int_0^\infty\frac{e^{-ay}}{1+y^2}\,dy
  -
  (\sin a)\int_0^\infty\frac{ye^{-ay}}{1+y^2}\,dy.
\]
Consequently, for $a\ge1$,
\[
\begin{aligned}
  \left|\int_0^a \frac{\sin x}{x}\,dx-\frac{\pi}{2}\right|
  &\le
  \int_0^\infty e^{-ay}\,dy
  +
  \int_0^\infty ye^{-ay}\,dy\\
  &=
  \frac1a+\frac1{a^2}
  \le
  \frac{2}{a}.
\end{aligned}
\]
\emph{Note.} The extracted exercise text displayed $\arctan(a)$ in the
place where the computation over the stated strip gives $\pi/2$.  With
the stated strip $0<x<a$, $0<y<\infty$, the formula above is the standard
Fubini computation.  Since $|\pi/2-\arctan(a)|\le 1/a$, it also gives
the related bound
\[
  \left|\int_0^a \frac{\sin x}{x}\,dx-\arctan(a)\right|\le \frac{3}{a}
  \qquad (a\ge1).
\]

\end{document}
